\documentclass[10pt,a4paper]{amsart}
\usepackage[margin=19mm]{geometry}
\usepackage{amsmath,amssymb,mathtools}
\usepackage[scr=rsfs,cal=euler]{mathalfa}
\usepackage{xfrac}
\usepackage[unicode,hidelinks,bookmarks=false]{hyperref}
\newcommand{\ie}{i.e.\ }
\newcommand{\cf}{cf.\ }
\newcommand{\resp}{respectively\ }
\newcommand{\Subsection}{\subsection}
\providecommand{\nobreakdash}{\nobreak\hbox{-}}
\setlength{\emergencystretch}{2em}
\multlinegap0pt
\usepackage{amssymb}
\usepackage{enumerate}
\usepackage[shortlabels]{enumitem}
\usepackage{mathtools}
\usepackage{xcolor}
\usepackage{dsfont}
\newtheorem{lemma}{Lemma}
\newcommand{\C}{\mathbb{C}}
\DeclareMathOperator{\Gal}{Gal}
\newcommand{\Q}{\mathbb{Q}}
\title{An exploration of degeneracy in abelian varieties of Fermat type}
\author{Heidi Goodson}
\date{Source: arXiv:2211.03909v2 (CC BY 4.0)}
\begin{document}
\maketitle

\noindent\emph{Source and licence.} An exploration of degeneracy in abelian varieties of Fermat type. Version arXiv:2211.03909v2 (CC BY 4.0), distributed by arXiv under the Creative Commons Attribution 4.0 International licence, http://creativecommons.org/licenses/by/4.0/, as recorded in the arXiv licence metadata for this exact version and shown on the abstract page https://arxiv.org/abs/2211.03909v2. Full licence notice: \texttt{licenses/arxiv-2211.03909-CC-BY-4.0.txt}.\par\smallskip

\noindent\textit{Editorial scope.} Counted unit: the article's lemma lemma:CMtype (source0.tex lines 497--501). The article's complete original proof of it is delivered with it (source0.tex lines 503--518, one \texttt{proof} environment of 16 lines). No mathematics is written, repaired or re-derived: the statement and the proof are the article's own lines, in the article's own notation, and no retained line is substituted or re-flowed.  Retained uncounted as the source-local context the proof needs: lines 477--479; lines 485--487.  Nothing was omitted from any retained span, so the package carries no omission record.  No retained line contains a citation, so the wrapper carries no bibliography and no bibliography entry.  No retained line contains a figure environment, an image-inclusion command or drawing code, so no image is delivered and the package declares no asset dependency.  Editorial formatting only, in addition: an AMS \texttt{amsart} preamble that restates the article's own declarations and macros used by the retained lines, namely the environments \texttt{lemma} on separate counters, together with the standard packages those lines need.  The retained environments print in this document as Lemma 1 in order of appearance, whereas the article prints the same items with the numbering of its own full document.  The complete native source of the article as distributed in the arXiv e-print for this exact version is bundled unchanged as \texttt{sources/133147/source0.tex}.
\begin{align}\label{eqn:Jpp_splitting}
    J_{p^2}\sim X\times J_p.
\end{align}

\begin{align}\label{eqn:Jpq_splitting}
    J_{pq}\sim X\times J_p \times J_q,
\end{align}

\begin{lemma}\label{lemma:CMtype}
Let $m=pq$, where $p$ and $q$ are odd primes (not necessarily distinct). Let $X$ be the Prym variety obtained as a quotient of $J_m$ as in \eqref{eqn:Jpp_splitting} or \eqref{eqn:Jpq_splitting}. Then the CM-type of $X$ is $(\Q(\zeta_m); \Phi)$, where
$$\Phi=\{\sigma_j \mid \gcd(j,m)=1\text{ and } j\leq g\},$$ 
$g=(m-1)/2$, and $\sigma_j(\zeta_m)=\zeta_m^j$.
\end{lemma}

\begin{proof}
For ease of notation, let $C\colon=C_m$. Since $C$ is a genus $g=(m-1)/2$ hyperelliptic curve, we may identify the complex uniformization of $J_m$ with $H^0(C_\C, \Omega_C^1)=\C\langle x^idx/y\mid 0\leq i< g\rangle$. We briefly split into two cases:

\textbf{Case 1: $p=q$.} Recall that there is a map 
 $\pi_p: (x,y)\mapsto (x^{p},y)$ from $C$ to $C_p$.  By pulling back regular differentials of $C_p$, we see that $\pi_p^*C_p$ corresponds to the vector subspace $$\C\left\langle{x^{ap-1}}dx/y\mid a=1, 2, \ldots,\frac{p-1}{2}\right\rangle$$ of $H^0(C_\C, \Omega_C^1)$. 
 Thus, the quotient $J_m/J_p$ (which is isogenous to $X$) has natural analytic uniformization given by the quotient of $\C\langle x^idx/y\mid 0\leq i< g\rangle$ by this vector space. The dimension of the new vector space is $\frac{m-1}{2}-\frac{p-1}{2}=\phi(m)/2$.

\textbf{Case 2: $p\not=q$.} In this case, there are two such maps:  $\pi_p: (x,y)\mapsto (x^{q},y)$ from $C$ to $C_p$ and $\pi_q: (x,y)\mapsto (x^{p},y)$ from $C$ to $C_q$. As above, we pull back regular differentials of $C_p$ and $C_q$. The quotient of $J_m$ by the Jacobians of the two curves (which is isogenous to $X$) has  natural analytic uniformization given by the quotient of
$\C\langle x^idx/y\mid 0\leq i< g\rangle$ by the vector subspace generated by
$$\C\left\langle{x^{aq-1}}dx/y\mid a=1, 2, \ldots,({p-1})/{2}\right\rangle \cup \C\left\langle{x^{ap-1}}dx/y\mid a=1, 2, \ldots,({q-1})/{2}\right\rangle.$$
The dimension of the new vector space is $\frac{m-1}{2}-\frac{p-1}{2}-\frac{q-1}{2}=\phi(m)/2$.

In both cases, the automorphism $\alpha\colon(x,y)\mapsto (\zeta_mx,y)$ induces an action of $\Q(\zeta_m)$ on the $(\phi(m)/2)$-dimensional vector space given above. Since $\alpha^*x^jdx/y=\zeta_m^{j+1} x^jdx/y$, we find that characters appearing in this uniformization will be of the form $\zeta_m\mapsto \zeta_m^j$ where $j$ is coprime to $m$. 

Thus, we can identify the CM type $\Phi$ of $\Gal(\Q(\zeta_m)/\Q)$ from the statement of the lemma as the CM-type of  $X$.
\end{proof}

\end{document}
